\section{Capacity Bounds}
\label{sec:capacity}

We now establish how many objects can be encoded in superposition as a function of dimension and sparsity.

\subsection{Sparsity Measure}

\begin{definition}[Co-occurrence Sparsity]
\label{def:sparsity}
For co-occurrence matrix $P \in [0,1]^{n \times n}$, define the \emph{sparsity} as:
\begin{equation}
S(P) = 1 - \frac{1}{n(n-1)} \sum_{i \neq j} P_{ij}
\end{equation}
This measures the fraction of pairs that rarely co-occur. $S = 1$ means no pairs co-occur; $S = 0$ means all pairs always co-occur.
\end{definition}

\begin{example}[Sparsity in Both Settings]
\label{ex:sparsity}
\begin{itemize}
    \item \textbf{Representational:} With $P_{ij} = p_i p_j$ and uniform $p_i = p$, we have $S = 1 - p^2$. For $p = 0.1$, sparsity is $S = 0.99$.
    \item \textbf{Reasoning:} For a binary tree with depth $D$, there are $2^D$ leaf traces, but only $O(D)$ pairs share any edge. Thus $S \approx 1 - O(D/2^D) \to 1$ as $D \to \infty$.
\end{itemize}
\end{example}

\begin{remark}[Global vs.\ Local Sparsity]
\label{rem:global_local}
A subtle distinction arises between \emph{global} and \emph{local} sparsity in reasoning:
\begin{itemize}
    \item \textbf{Global $P$:} The co-occurrence matrix over \emph{all possible} traces. This is sparse because most trace pairs never both lead to valid solutions.
    \item \textbf{Local $P^{(t)}$:} The co-occurrence matrix over traces \emph{active at timestep $t$}. During BFS-like exploration, the ``frontier'' of active traces can grow, making $P^{(t)}$ locally dense.
\end{itemize}

This apparent tension resolves as follows: even when the frontier is large, the \emph{relevant} traces (those contributing to the eventual answer) remain sparse. The model maintains superposition over the full frontier, but the interference that \emph{matters}---between traces that both reach the goal---remains controlled by global sparsity.

Formally, let $F_t$ denote the frontier at step $t$ with $|F_t| = O(b^t)$ for branching factor $b$. The local interference is:
\begin{equation}
\mathcal{L}_{\text{int}}^{(t)} = \sum_{i, j \in F_t} P^{(t)}_{ij} \langle \mathbf{v}_i, \mathbf{v}_j \rangle^2
\end{equation}
Even if $|F_t|$ is large, the capacity bound (Theorem~\ref{thm:capacity}) applies with dimension $d$: superposition succeeds when $|F_t| \leq O(d^2 \cdot S^{(t)} / \tau)$. The model handles frontier growth by either (a) having sufficient dimension $d$, or (b) using iterative refinement (Section~\ref{sec:timespace}) to prune the frontier over time.
\end{remark}

\subsection{Main Capacity Result}

\begin{theorem}[Capacity-Sparsity Bound]
\label{thm:capacity}
Let $\mathbf{W}^* = \arg\min_{\mathbf{W}} \mathcal{L}(\mathbf{W}, P)$ be an optimal encoding in $d$ dimensions. Define the \emph{effective number of encoded objects} as:
\begin{equation}
n_{\text{eff}} = |\{i : \|\mathbf{w}^*_i\| \geq 1/2\}|
\end{equation}
Then for any interference tolerance $\tau > 0$:\footnote{We use $\tau$ (tolerance) for the capacity constraint budget to distinguish from $\epsilon$ used elsewhere for approximation error bounds.}
\begin{equation}
n_{\text{eff}} \leq d + \frac{d(d-1)}{2} \cdot \frac{S(P)}{\tau}
\label{eq:capacity_bound}
\end{equation}
where $S(P)$ is the co-occurrence sparsity.
\end{theorem}

\begin{proof}
Let $\mathbf{G} = \mathbf{W}^\top \mathbf{W} \in \R^{n \times n}$ be the Gram matrix, with $G_{ij} = \langle \mathbf{w}_i, \mathbf{w}_j \rangle$.

\textbf{Step 1: Counting constraints.}
In $d$ dimensions, at most $d$ vectors can be mutually orthogonal. For $n_{\text{eff}} > d$ objects with $\|\mathbf{w}_i\| \geq 1/2$, there must exist pairs with $G_{ij} \neq 0$.

\textbf{Step 2: Interference budget.}
The interference loss is:
\begin{equation}
\mathcal{L}_{\text{int}} = \sum_{i \neq j} P_{ij} G_{ij}^2 \leq \tau
\end{equation}

By Cauchy-Schwarz, for unit-norm vectors $|G_{ij}| \leq 1$. For vectors with $\|\mathbf{w}_i\| \geq 1/2$, we have $|G_{ij}| \leq \|\mathbf{w}_i\| \|\mathbf{w}_j\| \leq 1$.

\textbf{Step 3: Packing argument.}
Consider the $\binom{n_{\text{eff}}}{2}$ pairs of effectively encoded objects. Each pair $(i,j)$ with $P_{ij} > 0$ contributes at least $P_{ij} G_{ij}^2$ to the interference loss.

The number of pairs with $P_{ij} > 0$ is at most $(1 - S) \cdot n_{\text{eff}}(n_{\text{eff}}-1)/2$.

For the interference constraint to be satisfied, we need:
\begin{equation}
\sum_{i \neq j} P_{ij} G_{ij}^2 \leq \tau
\end{equation}

\textbf{Step 4: Welch bound packing.}
Beyond the first $d$ orthogonal vectors, each additional vector must have non-zero inner product with at least one existing vector. The \emph{Welch bound} \citep{welch1974lower} provides a fundamental limit: for $n$ unit vectors in $\R^d$, the maximum coherence satisfies
\begin{equation}
\mu_{\max} := \max_{i \neq j} |G_{ij}| \geq \sqrt{\frac{n - d}{d(n-1)}}
\end{equation}
This implies that packing $n \gg d$ vectors requires $\mu_{\max} = \Omega(1/\sqrt{d})$, so the number of nearly-orthogonal directions (with $|G_{ij}| \leq \delta$) scales as $O(d^2/\delta^2)$.

Specifically, if we allow $|G_{ij}|^2 \leq \delta$ for pairs with $P_{ij} = 1-S$ (the ``dense'' pairs), the maximum packing is:
\begin{equation}
n_{\text{eff}} \leq d + \frac{d(d-1)}{2} \cdot \frac{1}{\delta}
\end{equation}

\textbf{Step 5: Combining with sparsity.}
The effective interference per dense pair is $P_{ij} G_{ij}^2 \leq (1-S) \cdot \delta$. To achieve total interference $\leq \tau$:
\begin{equation}
(1-S) \cdot \delta \cdot \frac{n_{\text{eff}}(n_{\text{eff}}-1)}{2} \leq \tau
\end{equation}

Solving for $n_{\text{eff}}$ with the packing constraint yields:
\begin{equation}
n_{\text{eff}} \leq d + \frac{d(d-1)}{2} \cdot \frac{S}{\tau}
\end{equation}
as claimed.
\end{proof}

\begin{corollary}[Linear Scaling with Sparsity]
\label{cor:linear}
For fixed dimension $d$ and interference tolerance $\tau$, the capacity scales linearly with sparsity:
\begin{equation}
n_{\text{eff}} = O(d^2 \cdot S / \tau)
\end{equation}
\end{corollary}

\begin{corollary}[Equivalence of Scaling]
\label{cor:equiv_scaling}
Both representational and reasoning superposition exhibit the same capacity scaling:
\begin{itemize}
    \item \textbf{Representational:} With activation probability $p$, sparsity is $S = 1 - p^2$, giving $n_{\text{eff}} = O(d^2 / p^2)$ for small $p$.
    \item \textbf{Reasoning:} With $K$ total traces and $k$ valid traces, if validity is sparse ($k \ll K$), then $S \approx 1$ and $n_{\text{eff}} = O(d^2)$.
\end{itemize}
The identical scaling confirms the phenomena are mathematically equivalent.
\end{corollary}

\subsection{Tightness of the Bound}

\begin{proposition}[Lower Bound]
\label{prop:lower}
There exist co-occurrence matrices $P$ and encodings $\mathbf{W}$ achieving:
\begin{equation}
n_{\text{eff}} \geq d + c \cdot d^2 \cdot S / \tau
\end{equation}
for a universal constant $c > 0$.
\end{proposition}

\begin{proof}[Proof Sketch]
Consider the construction of \citet{elhage2022toy} using regular polytopes. For $d = 2$ dimensions, $n = 2k$ features can be arranged as vertices of a regular $k$-gon and its reflection, achieving $\langle \mathbf{w}_i, \mathbf{w}_j \rangle = \cos(2\pi |i-j|/k)$.

For sparse features with $P_{ij} = p^2$, the interference is:
\begin{equation}
\mathcal{L}_{\text{int}} = p^2 \sum_{i \neq j} \cos^2(2\pi|i-j|/k) = O(p^2 k)
\end{equation}

Setting this equal to $\tau$ and solving: $k = O(\tau / p^2) = O(\tau \cdot d^2 \cdot S)$ for $d = 2$.

The construction generalizes to higher dimensions using cross-polytopes and their duals.
\end{proof}
